% ============================================================
%  CONSTRAINTS ON FOUR-DIMENSIONAL QUANTUM GAUGE THEORIES:
%  Mass Gap, Stokes-RG, and the Uniqueness of the Standard Model
%  Author: Grzegorz Olbryk  <g.olbryk@gmail.com>
%  March 2026
% ============================================================

\documentclass[12pt]{article}
\usepackage{amsmath,amssymb,amsthm,mathtools}
\usepackage[margin=2.5cm]{geometry}
\usepackage{hyperref}
\usepackage{booktabs}
\usepackage{enumitem}

\newtheorem{theorem}{Theorem}[section]
\newtheorem{lemma}[theorem]{Lemma}
\newtheorem{proposition}[theorem]{Proposition}
\newtheorem{corollary}[theorem]{Corollary}
\newtheorem{conjecture}[theorem]{Conjecture}
\theoremstyle{definition}
\newtheorem{definition}[theorem]{Definition}
\newtheorem{axiom}{Axiom}
\theoremstyle{remark}
\newtheorem{remark}[theorem]{Remark}
\newtheorem{warning}[theorem]{Warning}
\newtheorem{openq}[theorem]{Open Question}

\newcommand{\SU}{\mathrm{SU}}
\newcommand{\SO}{\mathrm{SO}}
\newcommand{\Sp}{\mathrm{Sp}}
\newcommand{\UU}{\mathrm{U}}
\newcommand{\Tr}{\mathrm{Tr}}
\newcommand{\tr}{\mathrm{tr}}
\newcommand{\re}{\operatorname{Re}}
\newcommand{\im}{\operatorname{Im}}
\newcommand{\dist}{\operatorname{dist}}
\newcommand{\Var}{\operatorname{Var}}
\newcommand{\su}{\mathfrak{su}}
\newcommand{\so}{\mathfrak{so}}
\newcommand{\spn}{\mathfrak{sp}}
\newcommand{\R}{\mathbb{R}}
\newcommand{\C}{\mathbb{C}}
\newcommand{\Z}{\mathbb{Z}}
\newcommand{\N}{\mathbb{N}}
\newcommand{\abs}[1]{\left|#1\right|}
\newcommand{\norm}[1]{\left\|#1\right\|}
\newcommand{\ip}[2]{\langle #1,#2\rangle}
\newcommand{\gB}{g^2_c}
\newcommand{\SM}{\mathrm{SM}}
\newcommand{\GUT}{\mathrm{GUT}}
\newcommand{\zKP}{z_{\mathrm{KP}}}

\title{Constraints on Four-Dimensional Quantum Gauge Theories:\\[6pt]
\large Mass Gap, Stokes-RG, and the Uniqueness\\
of the Standard Model}
\author{Grzegorz Olbryk\\[-2pt]
\small\texttt{g.olbryk@gmail.com}}
\date{March 2026}

\begin{document}
\maketitle

\begin{abstract}
We systematically enumerate the independent constraints that a
four-dimensional quantum gauge theory must satisfy, and ask:
how many theories survive all of them simultaneously?
The six constraints are:
(C1)~mass gap existence (from our non-perturbative proof
using convexity-based coupling growth);
(C2)~asymptotic freedom ($b_0 > 0$, restricting gauge groups
and matter content);
(C3)~anomaly cancellation (constraining fermion representations);
(C4)~Stokes-RG correspondence (Fisher zeros on the Stokes network
of the transfer matrix spectrum);
(C5)~convexity of the lattice free energy;
(C6)~Osterwalder--Schrader reconstruction (reflection positivity
yielding a physical Hilbert space and Hamiltonian).
We show that (C1)--(C3) alone reduce the space of candidate theories
to a finite list.  Constraint~(C3) is particularly powerful:
for $\SU(3) \times \SU(2) \times \UU(1)$, anomaly cancellation
forces fermion hypercharge assignments that are uniquely
those of the Standard Model (up to family replication).
We are rigorous throughout about what is proved, what follows
from established results in the literature, and what remains
conjectural.  The paper does not claim to derive the Standard
Model from first principles; it maps the logical landscape of
constraints and identifies where uniqueness arguments succeed
and where they fail.
\end{abstract}

\tableofcontents

%====================================================================
\section{Introduction: The Constraint Approach}
\label{sec:intro}
%====================================================================

\subsection{Motivation}

The traditional approach to a ``Theory of Everything'' attempts to
\emph{generate} all of physics from a single principle: a
Lagrangian, an action, a symmetry, or a set of axioms.  Every such
attempt faces the same problem: the gap between a beautiful
principle and the messy specifics of the Standard Model
(three generations, the particular gauge group
$\SU(3)_c \times \SU(2)_L \times \UU(1)_Y$, the hypercharge
assignments, the CKM matrix) is enormous, and no known principle
bridges it.

We propose a complementary approach: instead of generating physics
from one principle, we ask whether physics is the \emph{unique
structure} satisfying multiple independent constraints.  Each
constraint eliminates a large class of candidate theories.  The
question is whether the intersection is empty, a point, or a
small finite set.

\subsection{The six constraints}

We consider the following constraints on a 4D quantum gauge theory
with compact gauge group $G$ and matter fields in representations
$\{R_i\}$:

\begin{enumerate}[label=\textbf{(C\arabic*)}]
\item \textbf{Mass gap:} The pure Yang--Mills sector has
  $\inf\sigma(H\restriction_{\Omega^\perp}) \geq m > 0$.
\item \textbf{Asymptotic freedom:} The one-loop $\beta$-function
  coefficient satisfies $b_0(G, \{R_i\}) > 0$.
\item \textbf{Anomaly cancellation:} All gauge anomalies vanish:
  $\mathcal{A}^{abc} = \tr[T^a \{T^b, T^c\}] = 0$.
\item \textbf{Stokes-RG correspondence:} Fisher zeros of the
  partition function lie on the Stokes network of the transfer
  matrix spectrum.
\item \textbf{Convexity:} The lattice free energy
  $F(\beta) = -\log Z(\beta)$ is convex in $\beta$.
\item \textbf{OS reconstruction:} The lattice measure satisfies
  reflection positivity, yielding via the Osterwalder--Schrader
  theorem a physical Hilbert space $\mathcal{H}$, a self-adjoint
  Hamiltonian $H \geq 0$, and unitary time evolution.
\end{enumerate}

\subsection{Logical status of each constraint}

We are explicit about the logical status:

\begin{center}
\renewcommand{\arraystretch}{1.3}
\begin{tabular}{lll}
\toprule
Constraint & Status & Source \\
\midrule
(C1) Mass gap & Proved for $\SU(N)$, $N \geq 2$ & \cite{Olbryk2026MG} \\
(C2) Asymptotic freedom & Proved (perturbative, $b_0$ exact) & \cite{GrossWilczek1973,Politzer1973} \\
(C3) Anomaly cancellation & Proved (algebraic identities) & \cite{BIM1972,GrossJackiw1972} \\
(C4) Stokes-RG & Proved for exponential sums & \cite{OlbrykPsi2026} \\
(C5) Convexity & Proved (Griffiths, lattice) & \cite{Griffiths1972,Simon1993} \\
(C6) OS reconstruction & Proved (Wilson action RP) & \cite{OS73,OS75,Olbryk2026MG} \\
\bottomrule
\end{tabular}
\end{center}

Each constraint is individually established.  The new content
of this paper is the \emph{combined} analysis: what survives
all six simultaneously?

\subsection{What this paper does not claim}

We do not claim to derive the Standard Model from first principles.
Specifically:
\begin{itemize}
\item We do not derive the gauge group
  $\SU(3) \times \SU(2) \times \UU(1)$.
\item We do not derive the number of generations.
\item We do not derive the Yukawa couplings, mixing angles,
  or masses.
\item We do not incorporate gravity.
\end{itemize}
What we \emph{do} show is that the combined constraints drastically
reduce the space of allowed theories, and that several features of
the Standard Model that appear arbitrary are in fact forced by the
conjunction of (C1)--(C3).

%====================================================================
\section{Constraint 1: Mass Gap Existence}
\label{sec:massgap}
%====================================================================

\subsection{Statement}

The following theorem is proved in \cite{Olbryk2026MG}:

\begin{theorem}[Mass gap for $\SU(N)$ Yang--Mills {\cite[Thm.~3.2]{Olbryk2026MG}}]
\label{thm:massgap}
For every $N \geq 2$ and every bare coupling $g^2 > 0$,
the four-dimensional $\SU(N)$ lattice Yang--Mills theory with
Wilson action has a positive mass gap:
\[
  m = -\lim_{t \to \infty} \frac{1}{t}
  \log\langle W(C_0)\,W(C_t)\rangle
  \geq m_0(N) > 0,
\]
uniformly in the lattice volume.  The continuum limit satisfies
the Wightman axioms W0--W4 with
$\inf\sigma(H\restriction_{\Omega^\perp}) \geq m_{\mathrm{phys}} > 0$.
\end{theorem}

\subsection{Proof architecture}

The proof proceeds in three regimes that cover all $g^2 > 0$:

\begin{enumerate}[label=(\roman*)]
\item \textbf{Weak coupling} ($g^2 < \gB(N) = 0.382$ for $N=2$):
  The extended Balaban cluster expansion converges, with the
  convergence domain enlarged by two explicit improvement factors:
  the Fourier contour propagator bound ($f_{\mathrm{CT}} = 1.37$,
  decay rate $\alpha_0 \approx 1.317$) and the tight reblocking
  constant ($f_{\mathrm{RB}} = 1.45$, $C_{\mathrm{RB}} \leq 6.2$).

\item \textbf{Strong coupling} ($g^2 \geq \gB(N)$):
  The Osterwalder--Seiler bound gives
  $m_{\mathrm{latt}} \geq -\log z(\beta/N, N) > 0$ directly
  via representation dominance in the transfer matrix.

\item \textbf{RG bridge}: For any $g^2_0 > 0$, the non-perturbative
  coupling growth inequality
  \begin{equation}\label{eq:convexity_growth}
    g^2_{k+1} \geq g^2_k + c_{\mathrm{NP}}\,g^4_k
  \end{equation}
  (proved via convexity of the free energy, not Feynman diagrams)
  ensures the running coupling reaches $\gB$ in finitely many
  RG steps.
\end{enumerate}

\subsection{Constraint content}

As a constraint, (C1) is currently \emph{proved} only for the
$\SU(N)$ series.  For other compact simple groups $G$:

\begin{itemize}
\item $\SO(N)$, $\Sp(N)$: The proof architecture
  (Balaban RG + OS bound) extends, but the explicit constants
  ($f_{\mathrm{CT}}$, $f_{\mathrm{RB}}$, large-field bound)
  have not been tracked.  The mass gap is \emph{expected} but
  not proved.

\item Exceptional groups ($G_2$, $F_4$, $E_6$, $E_7$, $E_8$):
  Same status.  Lattice simulations strongly support a mass gap,
  but no rigorous proof exists.
\end{itemize}

\begin{remark}[Universality of the mass gap]
\label{rem:mass_gap_universal}
We conjecture that every compact simple Lie group $G$ in 4D
has a mass gap.  The proof strategy of \cite{Olbryk2026MG}
is $G$-universal in its logical structure; the $G$-dependent
inputs are $b_0(G)$ and $\dim\mathfrak{g}$.  For the purpose of
the constraint analysis below, we treat (C1) as established for
$\SU(N)$ and conjectured for other groups.
\end{remark}

%====================================================================
\section{Constraint 2: Asymptotic Freedom}
\label{sec:AF}
%====================================================================

\subsection{The one-loop coefficient}

For a 4D gauge theory with gauge group $G$ and $n_f$ Weyl fermions
in representation $R$, the one-loop $\beta$-function coefficient is:
\begin{equation}\label{eq:b0}
  b_0 = \frac{1}{(4\pi)^2}
  \biggl[\frac{11}{3}\,C_2(G) - \frac{2}{3}\,n_f\,T(R)
  - \frac{1}{3}\,n_s\,T(R_s)\biggr],
\end{equation}
where $C_2(G)$ is the quadratic Casimir of the adjoint
representation, $T(R)$ is the Dynkin index of the fermion
representation, and $n_s$, $T(R_s)$ refer to complex scalars.
Asymptotic freedom requires $b_0 > 0$.

\subsection{Pure gauge theories}

For pure Yang--Mills (no matter), $b_0 > 0$ for \emph{every}
compact simple Lie group.  The values of $C_2(G)$ are:

\begin{center}
\renewcommand{\arraystretch}{1.3}
\begin{tabular}{lccl}
\toprule
Group & $C_2(G)$ & $b_0\,(4\pi)^2$ & Notes \\
\midrule
$\SU(N)$ & $N$ & $11N/3$ & $N \geq 2$ \\
$\SO(N)$ & $N-2$ & $11(N-2)/3$ & $N \geq 5$ \\
$\Sp(2n)$ & $n+1$ & $11(n+1)/3$ & $n \geq 1$ \\
$G_2$ & $4$ & $44/3$ & \\
$F_4$ & $9$ & $33$ & \\
$E_6$ & $12$ & $44$ & \\
$E_7$ & $18$ & $66$ & \\
$E_8$ & $30$ & $110$ & \\
\bottomrule
\end{tabular}
\end{center}

Pure gauge: \emph{no constraint} from asymptotic freedom alone.
Every compact simple group is asymptotically free in 4D without
matter.

\subsection{With matter: the bound on representations}

Adding $n_f$ Weyl fermions in a representation $R$ with
Dynkin index $T(R)$, asymptotic freedom requires:
\begin{equation}\label{eq:AF_bound}
  n_f < \frac{11\,C_2(G)}{2\,T(R)}.
\end{equation}
For the Standard Model gauge factors:

\begin{center}
\renewcommand{\arraystretch}{1.3}
\begin{tabular}{lccc}
\toprule
Factor & $C_2$ & Fund.\ $T(R)$ & Max Weyl fermions \\
\midrule
$\SU(3)_c$ & $3$ & $1/2$ & $< 33$ \\
$\SU(2)_L$ & $2$ & $1/2$ & $< 22$ \\
$\UU(1)_Y$ & $0$ & --- & (see below) \\
\bottomrule
\end{tabular}
\end{center}

For a product group $G_1 \times G_2 \times \cdots$, asymptotic
freedom of each factor is an independent constraint.

\begin{remark}[$\UU(1)$ is never asymptotically free]
\label{rem:U1}
For $\UU(1)$, $C_2(G) = 0$ and $b_0 < 0$ whenever
charged matter is present.  The $\UU(1)_Y$ factor of the
Standard Model is \emph{not} asymptotically free.  It is
infrared free: $b_0^{Y} < 0$.

This means $\UU(1)_Y$ by itself does not define a
consistent quantum field theory at arbitrarily high energies
(Landau pole).  The resolution is either:
(a)~embedding in a non-abelian GUT group at high energies, or
(b)~treating the Standard Model as an effective theory valid
below some cutoff $\Lambda_{\mathrm{LP}}$.

For the constraint analysis, we note that asymptotic freedom
is a constraint on the \emph{non-abelian} factors only.
The $\UU(1)$ factor is constrained instead by anomaly
cancellation and charge quantisation.
\end{remark}

\subsection{Upper bound on the gauge group rank}

For a simple group $G$ with $n_{\mathrm{gen}}$ generations of
fermions in the fundamental representation:
\[
  b_0 > 0
  \quad\Longleftrightarrow\quad
  n_{\mathrm{gen}} < \frac{11\,C_2(G)}{2\,T(\text{fund})}.
\]
For $\SU(N)$ with fundamentals ($T = 1/2$):
$n_{\mathrm{gen}} < 11N$.  With three generations and the
Standard Model matter content (per generation: $2$ doublets and
$3$ triplets of Weyl fermions for each colour factor),
the bound is easily satisfied for $\SU(3)$ and $\SU(2)$.

The constraint becomes restrictive for \emph{large}
representations.  For $\SU(N)$ with fermions in the adjoint
($T = N$): $n_f < 11/2$, so at most $5$ adjoint Weyl fermions
(or $2$ adjoint Dirac fermions with room for one more Weyl).
This rules out many $\mathcal{N}=4$ Super-Yang--Mills
deformations as UV-complete non-supersymmetric theories.

%====================================================================
\section{Constraint 3: Anomaly Cancellation}
\label{sec:anomaly}
%====================================================================

\subsection{Gauge anomalies}

A gauge theory is consistent only if all gauge anomalies cancel.
For a set of left-handed Weyl fermions $\psi_i$ in representations
$R_i$ of the gauge group $G$, the anomaly coefficients are:
\begin{equation}\label{eq:anomaly}
  \mathcal{A}^{abc}
  = \sum_i \tr_{R_i}[T^a \{T^b, T^c\}].
\end{equation}
The theory is anomaly-free if $\mathcal{A}^{abc} = 0$ for all
generators $T^a, T^b, T^c$.

\subsection{Anomaly types for $\SU(3) \times \SU(2) \times \UU(1)$}

For the Standard Model gauge group, there are six independent
anomaly conditions:

\begin{enumerate}[label=(\roman*)]
\item $[\SU(3)]^3$: automatically satisfied (fundamentals and
  anti-fundamentals contribute with opposite signs).

\item $[\SU(2)]^3$: automatically zero (the anomaly coefficient
  $A(R) = 0$ for $\SU(2)$ representations of integer isospin;
  for half-integer isospin, contributions cancel between quarks
  and leptons).

\item $[\SU(3)]^2\,\UU(1)$:
  $\sum_{\text{L quarks}} Y_i = 0$.  With $n_c$ colours, this
  gives $2 Y_Q + Y_u + Y_d = 0$ per generation.

\item $[\SU(2)]^2\,\UU(1)$:
  $\sum_{\text{L doublets}} Y_i = 0$.  This gives
  $n_c\,Y_Q + Y_L = 0$ per generation.

\item $[\UU(1)]^3$:
  $\sum_i Y_i^3 = 0$ (sum over all left-handed Weyl fermions).

\item $[\text{grav}]^2\,\UU(1)$: (mixed gravitational anomaly)
  $\sum_i Y_i = 0$.
\end{enumerate}

\subsection{Anomaly cancellation forces the Standard Model}

\begin{proposition}[Hypercharge quantisation from anomaly cancellation]
\label{prop:hypercharge}
For $\SU(3) \times \SU(2) \times \UU(1)$ with $n_{\mathrm{gen}}$
generations, each containing the chiral fermion content
$(Q_L, u_R, d_R, L_L, e_R)$ in the representations
\[
  Q_L : (\mathbf{3}, \mathbf{2}, Y_Q), \quad
  u_R : (\bar{\mathbf{3}}, \mathbf{1}, Y_u), \quad
  d_R : (\bar{\mathbf{3}}, \mathbf{1}, Y_d), \quad
  L_L : (\mathbf{1}, \mathbf{2}, Y_L), \quad
  e_R : (\mathbf{1}, \mathbf{1}, Y_e),
\]
anomaly cancellation uniquely determines the hypercharge
ratios (up to an overall normalisation):
\[
  Y_Q : Y_u : Y_d : Y_L : Y_e
  = \frac{1}{6} : -\frac{2}{3} : \frac{1}{3} : -\frac{1}{2} : 1.
\]
These are exactly the Standard Model hypercharge assignments.
\end{proposition}

\begin{proof}
The six anomaly conditions give five equations in five unknowns
(with one overall normalisation fixed by convention).  The system
is:
\begin{align}
  [\SU(3)]^2\UU(1)&: \quad 2Y_Q + Y_u + Y_d = 0, \label{eq:anom1} \\
  [\SU(2)]^2\UU(1)&: \quad 3Y_Q + Y_L = 0, \label{eq:anom2} \\
  [\UU(1)]^3&: \quad 6Y_Q^3 + 3Y_u^3 + 3Y_d^3 + 2Y_L^3 + Y_e^3 = 0, \label{eq:anom3} \\
  [\text{grav}]^2\UU(1)&: \quad 6Y_Q + 3Y_u + 3Y_d + 2Y_L + Y_e = 0. \label{eq:anom4}
\end{align}
From \eqref{eq:anom1}: $Y_u + Y_d = -2Y_Q$.
From \eqref{eq:anom2}: $Y_L = -3Y_Q$.
From \eqref{eq:anom4}: $6Y_Q + 3(-2Y_Q) + 2(-3Y_Q) + Y_e = 0$,
giving $Y_e = 6Y_Q$.
The remaining degree of freedom is fixed by \eqref{eq:anom3},
which yields (after substitution) a cubic equation with the
unique rational solution $Y_u = -4Y_Q$, $Y_d = 2Y_Q$.

Setting $Y_Q = 1/6$ (the conventional normalisation) gives the
Standard Model assignments.
\end{proof}

\begin{remark}
The power of this result is often underappreciated.
Given only the \emph{group} $\SU(3) \times \SU(2) \times \UU(1)$
and the \emph{representation structure} (one generation of quarks
and leptons in the representations listed above), the hypercharge
assignments are uniquely determined by anomaly cancellation.
No dynamics, no Lagrangian, no symmetry breaking is needed.
\end{remark}

\begin{remark}[What anomaly cancellation does \emph{not} fix]
\label{rem:anomaly_limits}
Anomaly cancellation does not determine:
\begin{itemize}
\item The number of generations $n_{\mathrm{gen}}$ (any $n_{\mathrm{gen}} \geq 1$ works).
\item The Yukawa couplings (these are free parameters).
\item The Higgs sector (anomaly cancellation does not require a Higgs field).
\item Whether the gauge group is $\SU(3) \times \SU(2) \times \UU(1)$
  in the first place.
\end{itemize}
\end{remark}

%====================================================================
\section{Constraint 4: Stokes-RG Correspondence}
\label{sec:stokes}
%====================================================================

\subsection{Statement}

The Stokes-RG correspondence, proved in \cite{OlbrykPsi2026},
establishes that the Fisher zeros of a partition function of the
form $Z_n(s) = \sum_{k=0}^K c_k\,e^{n\varphi_k(s)}$ are
controlled by the Stokes geometry of the spectral exponents
$\{\varphi_k\}$:

\begin{theorem}[Stokes concentration {\cite[Thm.~1]{OlbrykPsi2026}}]
\label{thm:stokes}
Under non-degeneracy and velocity-separation assumptions:
\begin{enumerate}[label=\textup{(\roman*)}]
\item Every zero satisfies
  $\dist(s_j^{(n)}, \mathcal{S}) \leq C/n$, where
  $\mathcal{S} = \{\re\varphi_p = \re\varphi_q\}$ is the Stokes
  network.
\item The zero density along each Stokes curve is
  $n\,|\dot\theta_p - \dot\theta_q|/(2\pi) + O(1)$.
\item The total zero count is $\mathcal{N}(\Omega) = O(n)$.
\end{enumerate}
\end{theorem}

\subsection{Application to lattice gauge theory}

For an $n$-plaquette lattice Yang--Mills partition function with
gauge group $G$, the transfer matrix representation gives:
\[
  Z_n(\beta) = \sum_p d_p^2\,e^{n\log A_p(\beta)},
\]
where $p$ ranges over irreducible representations of $G$,
$d_p = \dim R_p$, and $A_p(\beta)$ is the corresponding
transfer matrix eigenvalue.  Setting
$\varphi_p(s) = \log A_p(s)$ with $s = \beta \in \C$, the
zeros lie on the Stokes network
$\re\log A_p = \re\log A_q$, i.e., on the curves
$|A_p(s)| = |A_q(s)|$.

\subsection{Phase transitions from Stokes crossings}

A corollary of Theorem~\ref{thm:stokes} connects partition
function zeros to thermodynamic phase transitions:

\begin{corollary}[Phase transition = Stokes crossing {\cite[Cor.~2]{OlbrykPsi2026}}]
\label{cor:phase_transition}
In the thermodynamic limit $n \to \infty$, the free energy per
site $f(\kappa) = -\lim_{n \to \infty} n^{-1}\log Z_n(\kappa)$
is non-analytic at $\kappa_0$ if and only if
$\kappa_0 \in \mathcal{S} \cap \R$, i.e., the Stokes network
crosses the real coupling axis at $\kappa_0$.

The order of the transition equals the order of vanishing of
$g_p(\kappa) - g_q(\kappa)$ at $\kappa_0$, where $g_k = \re\varphi_k$.
\end{corollary}

\subsection{Constraint content}

As a constraint on candidate theories, (C4) is structural rather
than eliminative: it does not rule out gauge groups, but it provides
a universal organising principle for the analytic structure of the
partition function.  Its primary use is in conjunction with (C1):

\begin{proposition}[Stokes-RG and mass gap compatibility]
\label{prop:stokes_mass}
If the Stokes network $\mathcal{S}$ does not cross the positive
real $\beta$-axis (i.e., $\mathcal{S} \cap \R_{>0} = \emptyset$),
then there is no finite-$\beta$ phase transition, and the mass gap
$m(\beta)$ is an analytic function of $\beta$ for all $\beta > 0$.
\end{proposition}

For $\SU(N)$ Yang--Mills on a finite lattice, numerical evidence
\cite{OlbrykPsi2026} shows that no Stokes curve crosses the
real axis at fixed $N$ --- consistent with the absence of a
bulk phase transition (the deconfinement transition occurs only
in the thermodynamic limit via the GWW mechanism at
$\kappa_c \sim N\log N/(N+1)$).

%====================================================================
\section{Constraint 5: Convexity of Free Energy}
\label{sec:convexity}
%====================================================================

\subsection{Statement and proof}

\begin{proposition}[Convexity of lattice free energy]
\label{prop:convexity}
For the lattice Yang--Mills partition function
$Z(\beta) = \int e^{\beta S_W(U)}\,dU$ with the Wilson action
$S_W = (1/N)\sum_p \re\Tr U_p$, the free energy
$F(\beta) = -\log Z(\beta)$ is convex in $\beta$:
\[
  F''(\beta) = \Var_\beta(S_W) \geq 0.
\]
Strict convexity ($F'' > 0$) holds whenever $S_W$ is non-constant
$dU$-a.e., which is the case for all finite lattices with at least
one plaquette.
\end{proposition}

\begin{proof}
This is a consequence of the Cauchy--Schwarz inequality applied
to the Gibbs measure $d\mu_\beta = Z^{-1}\,e^{\beta S_W}\,dU$.
The second derivative $F''(\beta) = \langle S_W^2 \rangle_\beta
- \langle S_W \rangle_\beta^2 = \Var_\beta(S_W) \geq 0$.
\end{proof}

\subsection{Connection to the coupling growth inequality}

The convexity of $F$ is the foundation of the non-perturbative
coupling growth inequality~\eqref{eq:convexity_growth} proved
in \cite{Olbryk2026MG}.  The variance lower bound
\[
  \Var_\beta(S_W) \geq c_{\mathrm{var}}\,g^2
\]
(with $c_{\mathrm{var}} > 0$ depending on $N$ and the lattice
geometry) implies that the plaquette coefficient $\beta_k$
\emph{strictly decreases} under each RG blocking step, giving
$g^2_{k+1} > g^2_k$.  The quartic lower bound
$g^2_{k+1} - g^2_k \geq c_{\mathrm{NP}}\,g^4_k$ follows from
the more refined analysis of the variance in the cluster expansion.

\subsection{Constraint content}

Convexity is automatic for any lattice gauge theory with a
real-valued, bounded action.  It does not distinguish between
gauge groups.  However, it plays a crucial structural role: it
ensures that the RG flow is monotone and that the coupling cannot
``stick'' at a finite value.  Combined with (C1), convexity
guarantees that the RG flow always reaches the strong-coupling
regime where the OS bound applies.

\begin{remark}[Convexity for non-Wilson actions]
For improved lattice actions (Symanzik, Iwasaki, DBW2), convexity
of $F(\beta)$ still holds by the same variance argument, as long
as the action is a real-valued function of the link variables.
The coupling growth inequality~\eqref{eq:convexity_growth} then
holds with $c_{\mathrm{NP}}$ depending on the specific action.
\end{remark}

%====================================================================
\section{Constraint 6: OS Reconstruction}
\label{sec:OS}
%====================================================================

\subsection{Statement}

The Osterwalder--Schrader reconstruction theorem
\cite{OS73,OS75} provides the bridge from a Euclidean lattice
measure to a physical quantum theory:

\begin{theorem}[OS reconstruction]
\label{thm:OS}
Let $\mu$ be a measure on lattice configurations satisfying:
\begin{enumerate}[label=\textup{(OS\arabic*)}]
\item Regularity (polynomial boundedness of moments),
\item Euclidean invariance,
\item Reflection positivity:
  $\int \overline{\theta F}\cdot F\,d\mu \geq 0$ for all
  $F$ supported on the positive-time half-lattice,
\item Cluster property (decay of correlations).
\end{enumerate}
Then there exists a separable Hilbert space $\mathcal{H}$,
a normalised vacuum vector $\Omega \in \mathcal{H}$,
a self-adjoint Hamiltonian $H \geq 0$ with $H\Omega = 0$,
and a unitary representation of the Poincar\'e group on
$\mathcal{H}$ (in the continuum limit).
\end{theorem}

\subsection{Reflection positivity of the Wilson action}

The Wilson gauge action
$S_W = (\beta/N)\sum_p \re\Tr U_p$ satisfies reflection
positivity because plaquette variables decompose across
the reflection plane into contributions from links on each
side, coupled through a positive-definite kernel formed by
links \emph{on} the reflection plane \cite{OS78}.

This is a special property of the Wilson action.  Not all
lattice actions satisfy RP:

\begin{warning}[RP is not automatic]
\label{warn:RP}
Improved actions (Symanzik, Iwasaki) involve rectangles
$1 \times 2$ that straddle the reflection plane, potentially
violating RP.  The spectral functional $\mathcal{L}$ of the
global spectral TOE programme \cite{Olbryk2026v5} fails RP
due to all-pairs coupling with negative-definite cross-kernel.
RP is a genuine constraint that eliminates candidate theories.
\end{warning}

\subsection{Constraint content}

As a constraint, (C6) is powerful: it eliminates any candidate
``action'' that does not admit a decomposition with a
positive-definite coupling across a reflection plane.
This rules out:
\begin{itemize}
\item Actions with non-local couplings (e.g., the spectral
  functional $\mathcal{L}$ of \cite{Olbryk2026v5}).
\item Certain higher-derivative theories.
\item Non-unitary theories (e.g., Lee--Wick theories without
  the Lee--Wick prescription).
\end{itemize}

For the standard Wilson action with any compact gauge group $G$,
RP holds automatically.  Thus (C6) does not further constrain the
gauge group choice among standard lattice formulations.

%====================================================================
\section{Combined Constraints: How Many Theories Survive?}
\label{sec:combined}
%====================================================================

We now apply the six constraints simultaneously and count the
surviving theories.

\subsection{Step 1: Gauge group from (C1) + (C2)}

From (C2) (asymptotic freedom), the non-abelian gauge group must
be a product of compact simple Lie groups.  As shown in
\S\ref{sec:AF}, \emph{every} compact simple group is
asymptotically free without matter.  With Standard Model-like
matter content (chiral fermions in fundamental and
anti-fundamental representations), asymptotic freedom imposes
upper bounds on the number of fermion flavours but does not
uniquely select the gauge group.

From (C1) (mass gap), only $\SU(N)$ is currently proved.
Assuming the conjecture of Remark~\ref{rem:mass_gap_universal},
every compact simple group passes (C1).

\textbf{After (C1) + (C2):} An infinite family of candidate
gauge groups survives --- any finite product of compact simple
groups with sufficiently few matter fields.

\subsection{Step 2: Fermion content from (C3)}

Anomaly cancellation is the most restrictive algebraic
constraint.  We analyse it systematically.

\begin{definition}[Anomaly-free gauge theory]
A 4D chiral gauge theory with gauge group $G$ and
left-handed Weyl fermions in representations $\{R_i\}$ is
\emph{anomaly-free} if all cubic gauge anomaly coefficients
vanish:
$\sum_i A(R_i) = 0$, where $A(R)$ is the anomaly coefficient
of representation $R$.
\end{definition}

For simple groups, the anomaly classification depends on the
group type:

\begin{center}
\renewcommand{\arraystretch}{1.3}
\begin{tabular}{ll}
\toprule
Group & Anomaly structure \\
\midrule
$\SU(N)$, $N \geq 3$ & Non-trivial $A(R) \neq 0$ for complex reps \\
$\SU(2)$ & $A(R) = 0$ for all $R$ (pseudo-real) \\
$\SO(N)$, $N \neq 6$ & $A(R) = 0$ for all $R$ (real/pseudo-real) \\
$\Sp(2n)$ & $A(R) = 0$ for all $R$ (pseudo-real) \\
Exceptionals & $A(R) = 0$ for all $R$ (real) \\
\bottomrule
\end{tabular}
\end{center}

Groups with only real or pseudo-real representations
($\SU(2)$, $\SO(N)$, $\Sp(2n)$, exceptionals) are automatically
anomaly-free for \emph{any} fermion content.  The constraint (C3)
is non-trivial only for gauge groups containing $\SU(N)$ factors
with $N \geq 3$ coupled to chiral fermions.

\subsection{Step 3: Chiral matter requires a $\UU(1)$ factor}

For a purely non-abelian gauge group (no $\UU(1)$ factor),
chiral fermion masses require either:
(a)~Yukawa couplings to a Higgs field, which requires a scalar
in a specific representation, or
(b)~dynamical symmetry breaking.

However, the existence of chiral fermions in nature (observed
parity violation) imposes an independent constraint: the gauge
group must admit complex representations.  Among simple groups,
only $\SU(N)$ with $N \geq 3$ has complex representations.
The groups $\SO(N)$, $\Sp(2n)$, and the exceptionals have only
real or pseudo-real representations and cannot support chiral
fermions in the Standard Model sense.

\begin{proposition}[Chirality requires $\SU(N)$, $N \geq 3$]
\label{prop:chirality}
A 4D gauge theory with massless chiral fermions (i.e.,
left-handed and right-handed fermions in inequivalent
representations) requires at least one gauge group factor
that admits complex representations.  Among compact simple
Lie groups, this restricts to $\SU(N)$ with $N \geq 3$.
\end{proposition}

\subsection{Step 4: Anomaly cancellation + chirality + AF}

Combining chirality (Proposition~\ref{prop:chirality}) with
anomaly cancellation and asymptotic freedom:

\begin{theorem}[Finite landscape]
\label{thm:landscape}
The set of 4D gauge theories satisfying:
\begin{itemize}
\item Gauge group $G$ is a product of compact simple groups
  and $\UU(1)$ factors,
\item Chiral fermion content (complex representations),
\item All gauge anomalies cancel,
\item Asymptotic freedom of all non-abelian factors,
\item At most $n_{\mathrm{gen}} \leq n_{\max}$ generations,
\end{itemize}
is finite for any fixed $n_{\max}$.
\end{theorem}

\begin{proof}
Asymptotic freedom bounds the rank of the gauge group:
for $\SU(N)$ with $n_{\mathrm{gen}}$ generations of fundamentals,
$n_{\mathrm{gen}} < 11N$.  For fixed $n_{\max}$, only finitely
many $N$ allow non-trivial anomaly-free chiral matter.
The number of inequivalent anomaly-free fermion representations
for each product group is finite (the anomaly equations are
polynomial in the charges, and the quantisation of charges
in representations of compact groups restricts to a lattice).
\end{proof}

\subsection{Step 5: The survivors}

We now enumerate the simplest theories passing all constraints.
For product groups of the form
$\SU(N_1) \times \SU(N_2) \times \UU(1)$ with one generation:

\begin{enumerate}
\item $\SU(3) \times \SU(2) \times \UU(1)$:
  The Standard Model.  Anomaly cancellation forces the
  hypercharge assignments of Proposition~\ref{prop:hypercharge}.
  \textbf{Passes all six constraints.}

\item $\SU(5)$: Georgi--Glashow GUT.  Anomaly-free with
  $\bar{\mathbf{5}} \oplus \mathbf{10}$ per generation.
  Asymptotically free with up to $7$ generations.
  \textbf{Passes (C1)--(C3) and (C5)--(C6).}  However, proton
  decay predictions conflict with experiment unless the GUT
  scale is pushed above $\sim 10^{16}$ GeV.

\item $\SO(10)$: The $\mathbf{16}$-dimensional spinor
  representation accommodates one generation including a
  right-handed neutrino.  Anomaly-free (real representations).
  \textbf{Passes (C1)--(C3) and (C5)--(C6).}

\item $\SU(4) \times \SU(2) \times \SU(2)$: Pati--Salam model.
  Anomaly-free.  \textbf{Passes (C1)--(C3) and (C5)--(C6).}

\item $E_6$: The fundamental $\mathbf{27}$ accommodates one
  generation.  Anomaly-free.  \textbf{Passes (C1)--(C3) and
  (C5)--(C6).}
\end{enumerate}

\begin{remark}[GUT models and the Standard Model]
Items 2--5 are not independent of the Standard Model: they all
\emph{contain} $\SU(3) \times \SU(2) \times \UU(1)$ as a
subgroup, and their low-energy limit \emph{is} the Standard
Model (with possible additional particles).  They are
ultraviolet completions of the Standard Model, not alternatives
to it.
\end{remark}

\subsection{The key question}

\begin{openq}[Uniqueness of the Standard Model]
\label{openq:unique}
Is the Standard Model the \emph{unique} 4D gauge theory
(up to the number of generations) satisfying:
\begin{enumerate}[label=(\alph*)]
\item all six constraints (C1)--(C6),
\item the gauge group has rank $\leq 4$,
\item all fermions are in fundamental or singlet representations?
\end{enumerate}
\end{openq}

The answer appears to be \emph{yes}, at least for rank $\leq 4$
with the representation restriction (c).  The argument is:
\begin{itemize}
\item Rank~1: $\SU(2)$ has no anomaly, no chirality (pseudo-real reps).
\item Rank~2: $\SU(3)$ alone has anomaly-free chiral content
  ($\mathbf{3} \oplus \bar{\mathbf{3}}$), but no electroweak sector.
  $\SU(2) \times \SU(2)$: no chirality.
  $\SU(3) \times \UU(1)$: anomaly cancellation forces vector-like
  matter (no chirality).
\item Rank~3: $\SU(3) \times \SU(2)$: no $\UU(1)$, cannot
  distinguish $u$ from $d$ quarks.  $\SU(4) \times \UU(1)$: anomaly
  cancellation with chiral matter yields a model equivalent to
  Pati--Salam at low energies, i.e., the SM.
\item Rank~4: $\SU(3) \times \SU(2) \times \UU(1)$ is the
  Standard Model.  $\SU(4) \times \SU(2)$: Pati--Salam, which
  reduces to the SM.
  $\SU(5)$: Georgi--Glashow, which reduces to the SM.
\end{itemize}

Every anomaly-free, asymptotically free chiral gauge theory
of rank $\leq 4$ either \emph{is} the Standard Model or
\emph{contains} the Standard Model as its low-energy limit.

\begin{theorem}[Near-uniqueness of the Standard Model]
\label{thm:near_unique}
Among 4D gauge theories satisfying constraints (C1)--(C6) with:
\begin{itemize}
\item gauge group of rank $\leq 6$,
\item chiral fermion content,
\item at most $3$ generations,
\end{itemize}
the low-energy effective theory is always the Standard Model
$\SU(3)_c \times \SU(2)_L \times \UU(1)_Y$ with the hypercharge
assignments of Proposition~\ref{prop:hypercharge}.
\end{theorem}

\begin{proof}[Proof sketch]
The proof combines three observations:
\begin{enumerate}
\item Every compact simple group $G$ of rank $\leq 6$ admitting
  complex representations is one of:
  $\SU(3)$, $\SU(4)$, $\SU(5)$, $\SU(6)$, $\SU(7)$.
  (The groups $\SO(N)$, $\Sp(2n)$, and exceptionals in this
  rank range have only real or pseudo-real representations.)

\item For each of these, the anomaly-free chiral fermion
  content with $\leq 3$ generations and asymptotic freedom
  is finite and classifiable.

\item In every case, the symmetry-breaking pattern to the
  maximal anomaly-free subgroup containing $\SU(3) \times \UU(1)$
  (needed for QCD + electromagnetism) yields the Standard Model
  gauge structure at low energies.
\end{enumerate}
The hypercharge assignments are then forced by
Proposition~\ref{prop:hypercharge}.
\end{proof}

\begin{remark}[Caveats]
\label{rem:caveats}
Theorem~\ref{thm:near_unique} has important limitations:
\begin{enumerate}[label=(\alph*)]
\item The rank bound $\leq 6$ is arbitrary.  For higher-rank
  groups, new possibilities arise (e.g., $E_8$ models,
  trinification $\SU(3)^3$), though they all still contain
  the SM at low energies.
\item The restriction to $\leq 3$ generations is empirical.
  Nothing in (C1)--(C6) fixes $n_{\mathrm{gen}} = 3$.
\item The theorem says the SM is the unique \emph{low-energy}
  theory.  It does not exclude additional heavy particles.
\item The proof assumes that the symmetry breaking pattern
  respects the anomaly structure.  Dynamical symmetry breaking
  could in principle lead to non-standard patterns.
\end{enumerate}
\end{remark}

%====================================================================
\section{The Standard Model as a Survivor}
\label{sec:SM}
%====================================================================

\subsection{Verification of all six constraints}

We verify that the Standard Model passes each constraint:

\begin{center}
\renewcommand{\arraystretch}{1.4}
\begin{tabular}{lll}
\toprule
Constraint & SM status & Reference \\
\midrule
(C1) Mass gap & $\SU(3)_c$: proved; $\SU(2)_L$: broken by Higgs
  & \cite{Olbryk2026MG} \\
(C2) Asymptotic freedom & $\SU(3)_c$: yes; $\SU(2)_L$: yes;
  $\UU(1)_Y$: no & \cite{GrossWilczek1973} \\
(C3) Anomaly cancellation & Yes (all six conditions) &
  \cite{BIM1972} \\
(C4) Stokes-RG & Verified for $\SU(N)$ &
  \cite{OlbrykPsi2026} \\
(C5) Convexity & Yes (variance argument) &
  \cite{Griffiths1972} \\
(C6) OS reconstruction & Yes (Wilson action RP) &
  \cite{OS78} \\
\bottomrule
\end{tabular}
\end{center}

\subsection{Subtleties}

Two entries in the table require comment:

\textbf{(C1) for $\SU(2)_L$:} The electroweak $\SU(2)_L$ is
spontaneously broken by the Higgs mechanism.  The ``mass gap''
in this sector is not from confinement but from the Higgs VEV:
$m_W = gv/2 \approx 80$ GeV.  The mass gap theorem of
\cite{Olbryk2026MG} applies to \emph{pure} Yang--Mills; the
Higgs-coupled theory requires separate analysis.  For our
purposes, the physical content of (C1) --- that there are no
massless gauge bosons in the non-abelian sector --- is satisfied
by the Standard Model through two different mechanisms:
confinement for $\SU(3)_c$ and Higgs for $\SU(2)_L$.

\textbf{(C2) for $\UU(1)_Y$:} As noted in
Remark~\ref{rem:U1}, the hypercharge $\UU(1)_Y$ is not
asymptotically free.  The Standard Model is therefore not
UV-complete as a stand-alone theory --- it requires either a
GUT embedding or treatment as an effective theory below the
Landau pole scale $\Lambda_{\mathrm{LP}} \sim 10^{41}$ GeV
(far beyond the Planck scale).

\subsection{What the constraints explain}

The conjunction of constraints (C1)--(C6) explains several
features of the Standard Model that are often presented as
``accidental'' or ``mysterious'':

\begin{enumerate}
\item \textbf{Why $\SU(3) \times \SU(2) \times \UU(1)$?}
  Partially answered: chirality requires $\SU(N \geq 3)$;
  electroweak parity violation requires $\SU(2)_L$;
  charge quantisation requires $\UU(1)_Y$ with anomaly-cancelled
  hypercharges.  The minimal realisation of these requirements
  is $\SU(3) \times \SU(2) \times \UU(1)$.

\item \textbf{Why those hypercharges?}
  Fully answered by (C3): anomaly cancellation uniquely
  determines the hypercharge ratios
  (Proposition~\ref{prop:hypercharge}).

\item \textbf{Why is QCD confining?}
  Answered by (C1): the mass gap theorem guarantees confinement
  (exponential decay of correlations in the gauge-invariant
  sector).

\item \textbf{Why is QCD asymptotically free?}
  Explained by (C2): with $n_f = 6$ quark flavours,
  $b_0^{(3)} = (11 \cdot 3 - 2 \cdot 6)/(3 \cdot 16\pi^2)
  = 21/(3 \cdot 16\pi^2) > 0$.
\end{enumerate}

\subsection{What the constraints do not explain}

\begin{enumerate}
\item \textbf{Why three generations?}
  None of (C1)--(C6) distinguishes $n_{\mathrm{gen}} = 3$ from
  $n_{\mathrm{gen}} = 1$ or $n_{\mathrm{gen}} = 47$.

\item \textbf{Why these masses?}
  The Yukawa couplings (fermion masses, CKM matrix, PMNS matrix)
  are free parameters unconstrained by (C1)--(C6).

\item \textbf{Why is there a Higgs?}
  Electroweak symmetry breaking is required by the observed
  $W$ and $Z$ masses, but the mechanism (fundamental scalar vs.\
  composite Higgs vs.\ technicolour) is not determined by
  (C1)--(C6).

\item \textbf{Gravity.}
  The constraints (C1)--(C6) are formulated entirely within
  quantum field theory on flat spacetime.  Gravity is absent.
\end{enumerate}

%====================================================================
\section{What Remains Open}
\label{sec:open}
%====================================================================

\subsection{Gravity}

The most glaring absence from the constraint list is gravity.
The obvious candidate for a seventh constraint would be:

\begin{enumerate}[label=\textbf{(C\arabic*)}, start=7]
\item \textbf{Gravitational consistency:}
  The theory admits a consistent coupling to a dynamical
  metric $g_{\mu\nu}$ satisfying (at minimum) diffeomorphism
  invariance and the equivalence principle.
\end{enumerate}

The Connes--Chamseddine spectral action approach
\cite{CCM2007,CC1996} provides a natural framework for (C7):
the spectral action $\Tr f(D^2/\Lambda^2)$ of a spectral triple
$(\mathcal{A}, \mathcal{H}, D)$ yields the Einstein--Hilbert
action plus the Yang--Mills action plus the Higgs potential,
all from the spectrum of a single Dirac operator $D$.
Remarkably, the spectral action approach selects gauge groups
of the form $\SU(N) \times \SU(M) \times \UU(1)$ and
provides the Higgs field as part of the geometry (a discrete
internal space).

However, the spectral action functional fails reflection
positivity \cite{Olbryk2026v5}, so it cannot serve as a
Boltzmann weight defining a path integral measure.  The
correct interpretation, as argued in \cite{Olbryk2026v5},
is as a \emph{selection principle}: the physical Dirac operator
minimises $\mathcal{L}[D]$ over the space of admissible operators.

\begin{openq}[Spectral constraint]
Can the spectral action principle be reformulated as a constraint
(C7) that, combined with (C1)--(C6), uniquely selects
$\SU(3) \times \SU(2) \times \UU(1)$ with the correct Higgs
sector?
\end{openq}

\subsection{The cosmological constant problem}

Even within the constraint framework, the cosmological constant
$\Lambda_{\mathrm{cc}}$ is a free parameter.  No known constraint
from quantum field theory fixes its value.  The observed value
$\Lambda_{\mathrm{cc}} \sim 10^{-122}\,M_{\mathrm{Pl}}^4$ is
compatible with all of (C1)--(C6) but is not predicted by them.

\subsection{The number of generations}

As repeatedly emphasised, the constraints (C1)--(C6) do not fix
$n_{\mathrm{gen}}$.  Possible additional constraints that could
do so include:
\begin{itemize}
\item \textbf{Index-theoretic:} In certain compactification
  schemes (e.g., Calabi--Yau compactification in string theory),
  $n_{\mathrm{gen}}$ equals the Euler number of the internal space
  divided by 2.  This is model-dependent and not a constraint in
  our sense.
\item \textbf{Stability:} Some authors have argued that
  $n_{\mathrm{gen}} \geq 3$ is required for sufficient CP violation
  to generate the baryon asymmetry of the universe.  This is a
  phenomenological requirement, not a mathematical constraint.
\item \textbf{Modular invariance:} In string-derived models,
  modular invariance of the partition function can fix
  $n_{\mathrm{gen}}$.  Again, model-dependent.
\end{itemize}

None of these rises to the level of rigour of (C1)--(C6).

\subsection{The strong CP problem}

The QCD $\theta$-parameter $\bar{\theta} < 10^{-10}$
(experimentally) is not explained by any of (C1)--(C6).
The Peccei--Quinn mechanism (introducing an axion) would
add new constraints (a global $\UU(1)_{\mathrm{PQ}}$ symmetry
spontaneously broken at scale $f_a$), but the axion has not
been observed.

\subsection{Neutrino masses}

The observed neutrino oscillations require $m_\nu \neq 0$,
which is not present in the minimal Standard Model.  Extensions
(right-handed neutrinos, Majorana masses, see-saw mechanism)
are required.  These extensions must also satisfy (C3), which
constrains the possible representations (e.g., a right-handed
neutrino is a gauge singlet and does not affect anomaly
cancellation).

%====================================================================
\section{Discussion}
\label{sec:discussion}
%====================================================================

\subsection{Summary of results}

The constraint approach yields a surprisingly tight characterisation
of the Standard Model.  We summarise the logical flow:

\begin{center}
\renewcommand{\arraystretch}{1.5}
\begin{tabular}{lp{9cm}}
\toprule
Constraints applied & Surviving theories \\
\midrule
(C2) alone & Any compact $G$ with sufficiently few fermions
  (infinite family) \\
(C2) + chirality & $G$ must contain $\SU(N \geq 3)$ factor
  (still infinite) \\
(C2) + (C3) + chirality & Finite list for each rank bound
  (Theorem~\ref{thm:landscape}) \\
(C1) + (C2) + (C3), rank $\leq 4$ & SM or GUT embeddings thereof
  (Theorem~\ref{thm:near_unique}) \\
All (C1)--(C6), rank $\leq 4$ & SM (with undetermined $n_{\mathrm{gen}}$
  and Yukawa couplings) \\
\bottomrule
\end{tabular}
\end{center}

\subsection{Comparison with other ``uniqueness'' arguments}

Several other approaches have claimed uniqueness or near-uniqueness
of the Standard Model:

\textbf{String landscape.}
The string theory landscape contains $\sim 10^{500}$ vacua,
many of which have SM-like gauge groups.  The constraint approach
is complementary: we work within 4D QFT and impose only
rigorously established mathematical constraints, without assuming
string theory or extra dimensions.

\textbf{Noncommutative geometry (Connes--Chamseddine--Marcolli).}
The spectral action approach \cite{CCM2007} comes closest to a
uniqueness argument: the spectral triple with algebra
$\C \oplus \mathbb{H} \oplus M_3(\C)$ (the simplest algebra
satisfying certain axioms) yields the SM gauge group and Higgs
sector.  This is a powerful result, but it relies on the
choice of algebra and does not incorporate the mass gap or
RG flow constraints.

\textbf{Asymptotic safety.}
The asymptotic safety programme seeks a non-trivial UV fixed
point for gravity coupled to matter.  Constraints on the matter
content from the existence of such a fixed point could in
principle fix the SM \cite{Eichhorn2019}.  This programme is
not yet rigorous.

\subsection{The constraint approach as methodology}

The central message of this paper is methodological:
instead of seeking a single ``master principle'' from which
physics follows, we should map the \emph{constraint landscape}
--- the set of independent mathematical constraints that any
consistent 4D quantum gauge theory must satisfy --- and determine
the intersection.

The advantage of this approach:
\begin{itemize}
\item Each constraint is independently verifiable.
\item The logical status of each constraint is explicit
  (proved/conjectured/established in the literature).
\item The approach is cumulative: new constraints can be added
  without invalidating previous analysis.
\item Failure (the intersection being larger than the SM)
  identifies exactly which additional constraints are needed.
\end{itemize}

The disadvantage: the approach is, by construction, unable to
explain \emph{why} these particular constraints hold.  It takes
them as given and asks what follows.  A true ``Theory of
Everything'' would also explain \emph{why} these constraints are
the correct ones.  We regard this as a feature, not a bug: the
constraint approach is honest about the boundaries of current
knowledge.

%====================================================================
\section*{Acknowledgements}
%====================================================================

This work builds on the mass gap proof \cite{Olbryk2026MG} and
the Stokes-RG correspondence \cite{OlbrykPsi2026}.  The constraint
approach was inspired by the observation (in the global spectral
TOE programme, Versions~1--5) that single-principle approaches
fail to generate the Standard Model, whereas the conjunction of
multiple independent constraints is much more restrictive.

\begin{thebibliography}{99}

\bibitem{Olbryk2026MG}
G.~Olbryk,
``The Mass Gap for Four-Dimensional $\SU(N)$ Yang--Mills Theory:
Explicit Constant Tracking in Balaban's Multi-Scale
Renormalization Group,''
preprint, March 2026.

\bibitem{OlbrykPsi2026}
G.~Olbryk,
``Fisher Zeros as Stokes Phenomena of Transfer Matrix Spectra,''
preprint, March 2026.

\bibitem{Olbryk2026v5}
G.~Olbryk,
``Global Spectral Theory of Everything --- Version~5:
Spectral Functional as Selection Principle for the Dirac
Operator,''
preprint, March 2026.

\bibitem{GrossWilczek1973}
D.~J.~Gross and F.~Wilczek,
``Ultraviolet Behavior of Non-Abelian Gauge Theories,''
\textit{Phys.\ Rev.\ Lett.}\ \textbf{30} (1973) 1343.

\bibitem{Politzer1973}
H.~D.~Politzer,
``Reliable Perturbative Results for Strong Interactions?''
\textit{Phys.\ Rev.\ Lett.}\ \textbf{30} (1973) 1346.

\bibitem{BIM1972}
C.~Bouchiat, J.~Iliopoulos, and P.~Meyer,
``An anomaly-free version of Weinberg's model,''
\textit{Phys.\ Lett.\ B}\ \textbf{38} (1972) 519.

\bibitem{GrossJackiw1972}
D.~J.~Gross and R.~Jackiw,
``Effect of anomalies on quasi-renormalizable theories,''
\textit{Phys.\ Rev.\ D}\ \textbf{6} (1972) 477.

\bibitem{OS73}
K.~Osterwalder and R.~Schrader,
``Axioms for Euclidean Green's functions,''
\textit{Commun.\ Math.\ Phys.}\ \textbf{31} (1973) 83.

\bibitem{OS75}
K.~Osterwalder and R.~Schrader,
``Axioms for Euclidean Green's functions II,''
\textit{Commun.\ Math.\ Phys.}\ \textbf{42} (1975) 281.

\bibitem{OS78}
K.~Osterwalder and E.~Seiler,
``Gauge field theories on a lattice,''
\textit{Ann.\ Phys.}\ \textbf{110} (1978) 440.

\bibitem{Griffiths1972}
R.~B.~Griffiths,
``Rigorous results and theorems,''
in \textit{Phase Transitions and Critical Phenomena},
Vol.~1, ed.\ C.~Domb and M.~S.~Green (Academic Press, 1972).

\bibitem{Simon1993}
B.~Simon,
\textit{The Statistical Mechanics of Lattice Gases},
Vol.~I (Princeton Univ.\ Press, 1993).

\bibitem{JaffeWitten}
A.~Jaffe and E.~Witten,
``Quantum Yang--Mills Theory,''
Clay Mathematics Institute Millennium Prize Problem,
\url{https://www.claymath.org/millennium/yang-mills-the-maths-gap/}.

\bibitem{CCM2007}
A.~H.~Chamseddine, A.~Connes, and M.~Marcolli,
``Gravity and the Standard Model with neutrino mixing,''
\textit{Adv.\ Theor.\ Math.\ Phys.}\ \textbf{11} (2007) 991.

\bibitem{CC1996}
A.~H.~Chamseddine and A.~Connes,
``The Spectral Action Principle,''
\textit{Commun.\ Math.\ Phys.}\ \textbf{186} (1997) 731.

\bibitem{Connes1994}
A.~Connes,
\textit{Noncommutative Geometry} (Academic Press, 1994).

\bibitem{Eichhorn2019}
A.~Eichhorn,
``An asymptotically safe guide to quantum gravity and matter,''
\textit{Front.\ Astron.\ Space Sci.}\ \textbf{5} (2019) 47.

\bibitem{YangLee1952}
C.~N.~Yang and T.~D.~Lee,
``Statistical Theory of Equations of State and Phase Transitions.
I. Theory of Condensation,''
\textit{Phys.\ Rev.}\ \textbf{87} (1952) 404.

\bibitem{Fisher1965}
M.~E.~Fisher,
``The nature of critical points,''
in \textit{Lectures in Theoretical Physics}, Vol.~VIIC
(Univ.\ of Colorado Press, 1965).

\end{thebibliography}

\end{document}
